\documentclass[aspectratio=169,10pt,english]{beamer}

\input{../../../_preambulo_beamer.tex}
\input{../../../_comandos_beamer.tex}

\selectlanguage{english}

\title{MA1001B - Statistical Modeling for Decision Making}
\subtitle{Lesson 13: Discrete Random Variables: PMF, CDF, Expectation, Variance}
\author{}
\institute{Tecnol\'ogico de Monterrey}
\date{}

\begin{document}

\begin{frame}[plain]
  \vspace{1.2cm}
  {\Large\bfseries MA1001B - Statistical Modeling for Decision Making\par}
  \vspace{0.7cm}
  {\large Lesson 13: Discrete Random Variables\par}
  \vspace{0.15cm}
  {\large PMF, CDF, Expectation, Variance\par}
  \vspace{0.7cm}
  {\color{TecRojo}\rule{\textwidth}{1.2pt}}
  \vspace{0.8cm}

  MA1001B Faculty\\[0.3cm]
  Term\\[0.3cm]
  Tecnol\'ogico de Monterrey
\end{frame}

\begin{frame}{A Defect Count as a Random Variable}
  \begin{alertblock}{Guiding question}
    How is a discrete random variable specified, and why is $E[X^2]$ not the
    same number as $(E[X])^2$?
  \end{alertblock}

  \vspace{0.6em}
  By the end of this lesson, you can:
  \begin{itemize}
    \item check that a PMF has nonnegative masses that sum to 1;
    \item accumulate the PMF into a CDF;
    \item compute $E[X]$ and $E[X^2]$ as weighted sums;
    \item obtain $\Var(X)=E[X^2]-(E[X])^2$.
  \end{itemize}

  \vfill
  \scriptsize All defect counts used today are synthetic. No real company data are used.
\end{frame}

\begin{frame}{Support and Probability Masses}
  \small
  Let $X$ be the defect count on a synthetic lot, $X\in\{0,1,2,3\}$.
  \[
    P(X=0)=0.50,\quad
    P(X=1)=0.30,\quad
    P(X=2)=0.15,\quad
    P(X=3)=0.05.
  \]
  A function $p$ is a PMF when $p(x)\ge 0$ for every $x$ and
  $\sum_x p(x)=1$.
\end{frame}

\begin{frame}[fragile]{Step 1: Probability Mass Function}
  \begin{columns}[T]
    \begin{column}{0.42\textwidth}
      \small
      \begin{lstlisting}
support = [0, 1, 2, 3]
pmf = [0.50, 0.30, 0.15, 0.05]
mass_sum = sum(pmf)
      \end{lstlisting}

      \begin{block}{Verified output}
        $P(X=0)=0.500000$\\
        $P(X=1)=0.300000$\\
        $P(X=2)=0.150000$\\
        Sum $=1.000000$
      \end{block}
    \end{column}
    \begin{column}{0.58\textwidth}
      \centering
      \includegraphics[width=\textwidth]{../figures/discrete_rv_01_pmf.png}
      \vspace{0.2em}
      \scriptsize PMF stem plot from Step 1
    \end{column}
  \end{columns}
\end{frame}

\begin{frame}{The CDF Accumulates the PMF}
  \small
  \[
    F(x)=P(X\le x)=\sum_{t\le x}P(X=t).
  \]
  Verified values:
  \[
    F(0)=0.50,\quad F(1)=0.80,\quad F(2)=0.95,\quad F(3)=1.00.
  \]
  Interval probabilities come from differences:
  \[
    P(1<X\le 3)=F(3)-F(1)=0.20.
  \]
\end{frame}

\begin{frame}[fragile]{Step 2: Cumulative Distribution Function}
  \begin{columns}[T]
    \begin{column}{0.40\textwidth}
      \small
      \begin{lstlisting}
cdf = np.cumsum(pmf)
p_hi = cdf[3] - cdf[1]
      \end{lstlisting}

      \begin{block}{Verified output}
        $F(1)=0.800000$\\
        $F(3)=1.000000$\\
        $P(1<X\le 3)=0.200000$
      \end{block}
    \end{column}
    \begin{column}{0.60\textwidth}
      \centering
      \includegraphics[width=\textwidth]{../figures/discrete_rv_02_cdf.png}
      \vspace{0.2em}
      \scriptsize Step-function CDF from Step 2
    \end{column}
  \end{columns}
\end{frame}

\begin{frame}{Step 3: $E[X^2]$ Is Not $(E[X])^2$}
  \begin{columns}[T]
    \begin{column}{0.42\textwidth}
      \small
      \[
        E[X]=0.750000,
      \]
      \[
        (E[X])^2=0.562500,
      \]
      \[
        E[X^2]=1.350000,
      \]
      \[
        \Var(X)=0.787500.
      \]

      \vspace{0.3em}
      \textbf{Limit:} $E[X^2]\ne(E[X])^2$. Lesson 14 specializes this language
      to Bernoulli trials and the binomial model.
    \end{column}
    \begin{column}{0.58\textwidth}
      \centering
      \includegraphics[width=\textwidth]{../figures/discrete_rv_03_expectation_variance.png}
      \vspace{0.2em}
      \scriptsize Moments and variance from Step 3
    \end{column}
  \end{columns}
\end{frame}

\begin{frame}{Concept Check}
  \begin{enumerate}
    \item Why is $P(X=4)=0$ for this $X$?
    \item Compute $F(2)$ from the PMF.
    \item Compute $E[X]$ from the four masses.
    \item Why is $1.350000$ not equal to $0.562500$?
  \end{enumerate}

  \vspace{0.6em}
  \begin{block}{Connection to Lesson 14}
    The session can continue directly into Bernoulli trials, the binomial
    PMF, and a constant-success-probability assumption.
  \end{block}

  \vfill
  \scriptsize Source: OpenStax, \textit{Introductory Statistics 2e},
  Sections 4.1--4.2. CC BY-NC-SA.
\end{frame}

\appendix



\end{document}
